\documentclass[10pt,a4paper]{amsart}
\usepackage{amsmath,amssymb,amsthm,graphicx}
% Retained preamble for the Stacks Project chapter snapshot in this directory.
% The chapter's own preamble.tex is a copy of the shared Stacks Project preamble; the block
% below is sliced MECHANICALLY, line for line, from it. Nothing was added, removed or
% redefined; the wrapper class, the xr-hyper/hyperref wiring and the externaldocument list
% are omitted because the assembler supplies the document class.
%
% lines taken from preamble.tex: 3, 5, 9, 11, 12, 156, 157, 158, ... (51 lines total)
%
% package coverage restored from the chapter's own preamble (2026-10-13): xy, xy-2cell

\usepackage[all]{xy}
\xyoption{2cell}
\UseAllTwocells

\newenvironment{reference}{\comment}{\endcomment}
\newenvironment{slogan}{\comment}{\endcomment}
\newenvironment{history}{\comment}{\endcomment}
\theoremstyle{plain}
\newtheorem{theorem}[subsection]{Theorem}
\newtheorem{proposition}[subsection]{Proposition}
\newtheorem{lemma}[subsection]{Lemma}
\theoremstyle{definition}
\newtheorem{definition}[subsection]{Definition}
\newtheorem{example}[subsection]{Example}
\newtheorem{exercise}[subsection]{Exercise}
\newtheorem{situation}[subsection]{Situation}
\theoremstyle{remark}
\newtheorem{remark}[subsection]{Remark}
\newtheorem{remarks}[subsection]{Remarks}
\numberwithin{equation}{subsection}
\def\lim{\mathop{\mathrm{lim}}\nolimits}
\def\colim{\mathop{\mathrm{colim}}\nolimits}
\def\Spec{\mathop{\mathrm{Spec}}}
\def\Hom{\mathop{\mathrm{Hom}}\nolimits}
\def\Ext{\mathop{\mathrm{Ext}}\nolimits}
\def\SheafHom{\mathop{\mathcal{H}\!\mathit{om}}\nolimits}
\def\SheafExt{\mathop{\mathcal{E}\!\mathit{xt}}\nolimits}
\def\Sch{\mathit{Sch}}
\def\Mor{\mathop{\mathrm{Mor}}\nolimits}
\def\Ob{\mathop{\mathrm{Ob}}\nolimits}
\def\Sh{\mathop{\mathit{Sh}}\nolimits}
\def\NL{\mathop{N\!L}\nolimits}
\def\CH{\mathop{\mathrm{CH}}\nolimits}
\def\proetale{{pro\text{-}\acute{e}tale}}
\def\etale{{\acute{e}tale}}
\def\QCoh{\mathit{QCoh}}
\def\Ker{\mathop{\mathrm{Ker}}}
\def\Im{\mathop{\mathrm{Im}}}
\def\Coker{\mathop{\mathrm{Coker}}}
\def\Coim{\mathop{\mathrm{Coim}}}
\DeclareMathSymbol{\boxtimes}{\mathbin}{AMSa}{"02}
\def\QCohstack{\mathcal{QC}\!\mathit{oh}}
\def\Cohstack{\mathcal{C}\!\mathit{oh}}
\def\Spacesstack{\mathcal{S}\!\mathit{paces}}
\def\Quotfunctor{\mathrm{Quot}}
\def\Hilbfunctor{\mathrm{Hilb}}
\def\Curvesstack{\mathcal{C}\!\mathit{urves}}
\def\Polarizedstack{\mathcal{P}\!\mathit{olarized}}
\def\Complexesstack{\mathcal{C}\!\mathit{omplexes}}
\def\Pic{\mathop{\mathrm{Pic}}\nolimits}
\def\Picardstack{\mathcal{P}\!\mathit{ic}}
\def\Picardfunctor{\mathrm{Pic}}
\def\Deformationcategory{\mathcal{D}\!\mathit{ef}}
\usepackage{cleveref}
\crefname{originallemma}{original lemma}{original lemmas}
\Crefname{originallemma}{Original lemma}{Original lemmas}
\makeatletter
\newlabel{bootstrap-lemma-descend-algebraic-space}{{0ADV}{source}{}{sourcepointer-1}{}}
\newlabel{bootstrap-lemma-descend-algebraic-space@cref}{{[originallemma][0][]0ADV}{[1][0][]source}{}{}{}}
\makeatother
\crefname{originallemma}{original lemma}{original lemmas}
\Crefname{originallemma}{Original lemma}{Original lemmas}
\makeatletter
\newlabel{lemma-pre-stack-of-spaces}{{04UA}{source}{}{sourcepointer-2}{}}
\newlabel{lemma-pre-stack-of-spaces@cref}{{[originallemma][0][]04UA}{[1][0][]source}{}{}{}}
\makeatother
\crefname{originallemma}{original lemma}{original lemmas}
\Crefname{originallemma}{Original lemma}{Original lemmas}
\makeatletter
\newlabel{lemma-spaces-strongly-cartesian}{{04U9}{source}{}{sourcepointer-3}{}}
\newlabel{lemma-spaces-strongly-cartesian@cref}{{[originallemma][0][]04U9}{[1][0][]source}{}{}{}}
\makeatother
\crefname{originallemma}{original lemma}{original lemmas}
\Crefname{originallemma}{Original lemma}{Original lemmas}
\makeatletter
\newlabel{spaces-descent-lemma-descending-property-locally-finite-presentation}{{041T}{source}{}{sourcepointer-4}{}}
\newlabel{spaces-descent-lemma-descending-property-locally-finite-presentation@cref}{{[originallemma][0][]041T}{[1][0][]source}{}{}{}}
\makeatother
\crefname{originallemma}{original lemma}{original lemmas}
\Crefname{originallemma}{Original lemma}{Original lemmas}
\makeatletter
\newlabel{spaces-morphisms-lemma-base-change-finite-type}{{03XH}{source}{}{sourcepointer-5}{}}
\newlabel{spaces-morphisms-lemma-base-change-finite-type@cref}{{[originallemma][0][]03XH}{[1][0][]source}{}{}{}}
\makeatother
\crefname{originallemma}{original lemma}{original lemmas}
\Crefname{originallemma}{Original lemma}{Original lemmas}
\makeatletter
\newlabel{stacks-definition-stack}{{026F}{source}{}{sourcepointer-6}{}}
\newlabel{stacks-definition-stack@cref}{{[originallemma][0][]026F}{[1][0][]source}{}{}{}}
\makeatother
\crefname{originallemma}{original lemma}{original lemmas}
\Crefname{originallemma}{Original lemma}{Original lemmas}
\makeatletter
\newlabel{stacks-lemma-presheaf-mor-map-fibred-categories}{{042V}{source}{}{sourcepointer-7}{}}
\newlabel{stacks-lemma-presheaf-mor-map-fibred-categories@cref}{{[originallemma][0][]042V}{[1][0][]source}{}{}{}}
\makeatother
\title{Finite-type algebraic spaces form a stack}
\author{The Stacks Project Authors}
\date{stacks-project:examples-stacks.tex:local-snapshot-20261012; GNU FDL 1.2 or later}
\begin{document}
\maketitle
\noindent\textit{Editorial note.} Statement span is lines 752-758 of sources/stacks-examples-stacks/pinned-original-source.tex; proof span is lines 760-823, the complete author proof (62 body lines).\par\medskip
\noindent\textit{Editorial note.} Navigation stubs: 7 cited-but-not-included labels are supplied as non-load-bearing pointers with their pinned tags where the snapshot tag index has them.\par\medskip
\noindent\textit{Editorial note.} extra\_packages is empty on purpose: the retained preamble already loads xy with options and a second declaration clashes in the assembler wrapper.\par\medskip
\noindent\textit{Editorial note.} Text is retained verbatim, including the author's own typos and placeholders.\par\medskip
\noindent\textit{Editorial note.} Queue index 1 of queue\_examples-stacks.json (pb=2629, stmt 752-758, proof 760-823, halfA).\par\medskip
\noindent\textit{Editorial note.} stacks-examples-stacks wave-1 (exs1): fresh adjudication accept (2026-10-14); pool per Chief 09:4x authorization (idle-line throughput over kun's backup note, SMC precedent); examples-stacks.tex 27 labels/20 with proof judged 9 accept / 11 reject; claim\_key v2 stacks:examples-stacks.tex:<label>.\par\medskip
\noindent\textit{Allowed standard prerequisites (editorial).}\begin{itemize}
\item Explicit prerequisite boundary. The cited inputs are prerequisites; the argument in the retained span is not replaced by any cited result.
\end{itemize}\medskip
\ifdefined\hypertarget\hypertarget{sourcepointer-1}{}\fi
\noindent\textit{Source reference.} Original lemma 0ADV: sources/stacks-examples-stacks/pinned-original-source.tex, source lines 823--823. Cited by the retained statement or proof but not part of the collected spans; supplied only as a navigation pointer. The argument does not depend on its content.\par\medskip
\ifdefined\hypertarget\hypertarget{sourcepointer-2}{}\fi
\noindent\textit{Source reference.} Original lemma 04UA: sources/stacks-examples-stacks/pinned-original-source.tex, source lines 675--679. Cited by the retained statement or proof but not part of the collected spans; supplied only as a navigation pointer. The argument does not depend on its content.\par\medskip
\ifdefined\hypertarget\hypertarget{sourcepointer-3}{}\fi
\noindent\textit{Source reference.} Original lemma 04U9: sources/stacks-examples-stacks/pinned-original-source.tex, source lines 633--637. Cited by the retained statement or proof but not part of the collected spans; supplied only as a navigation pointer. The argument does not depend on its content.\par\medskip
\ifdefined\hypertarget\hypertarget{sourcepointer-4}{}\fi
\noindent\textit{Source reference.} Original lemma 041T: sources/stacks-examples-stacks/pinned-original-source.tex, source lines 823--823. Cited by the retained statement or proof but not part of the collected spans; supplied only as a navigation pointer. The argument does not depend on its content.\par\medskip
\ifdefined\hypertarget\hypertarget{sourcepointer-5}{}\fi
\noindent\textit{Source reference.} Original lemma 03XH: sources/stacks-examples-stacks/pinned-original-source.tex, source lines 823--823. Cited by the retained statement or proof but not part of the collected spans; supplied only as a navigation pointer. The argument does not depend on its content.\par\medskip
\ifdefined\hypertarget\hypertarget{sourcepointer-6}{}\fi
\noindent\textit{Source reference.} Original lemma 026F: sources/stacks-examples-stacks/pinned-original-source.tex, source lines 823--823. Cited by the retained statement or proof but not part of the collected spans; supplied only as a navigation pointer. The argument does not depend on its content.\par\medskip
\ifdefined\hypertarget\hypertarget{sourcepointer-7}{}\fi
\noindent\textit{Source reference.} Original lemma 042V: sources/stacks-examples-stacks/pinned-original-source.tex, source lines 823--823. Cited by the retained statement or proof but not part of the collected spans; supplied only as a navigation pointer. The argument does not depend on its content.\par\medskip
\begin{lemma}
\label{lemma-stack-of-finite-type-spaces}
The functor
$p_{ft} : \Spacesstack_{ft} \to (\Sch/S)_{fppf}$
satisfies the conditions (1), (2) and (3) of
Stacks, Definition \ref{stacks-definition-stack}.
\end{lemma}

\begin{proof}
We are going to write this out in ridiculous detail (which may make
it hard to see what is going on).

\medskip\noindent
We have seen in
Lemma \ref{lemma-spaces-strongly-cartesian}
that a morphism $(f, g) : X/U \to Y/V$ of $\Spacesstack$ is
strongly cartesian if the induced morphism $f : X \to U \times_V Y$
is an isomorphism. Note that if $Y \to V$ is of finite type
then also $U \times_V Y \to U$ is of finite type, see
Morphisms of Spaces,
Lemma \ref{spaces-morphisms-lemma-base-change-finite-type}.
So if $(f, g) : X/U \to Y/V$ of $\Spacesstack$ is
strongly cartesian in $\Spacesstack$ and $Y/V$ is an object
of $\Spacesstack_{ft}$ then automatically also $X/U$ is an
object of $\Spacesstack_{ft}$, and of course $(f, g)$ is
also strongly cartesian in $\Spacesstack_{ft}$. In this way
we conclude that $\Spacesstack_{ft}$ is a fibred category over
$(\Sch/S)_{fppf}$. This proves (1).

\medskip\noindent
The argument above also shows that the inclusion
functor $\Spacesstack_{ft} \to \Spacesstack$ transforms
strongly cartesian morphisms into strongly cartesian morphisms.
In other words $\Spacesstack_{ft} \to \Spacesstack$ is
a $1$-morphism of fibred categories over $(\Sch/S)_{fppf}$.

\medskip\noindent
Let $U \in \Ob((\Sch/S)_{fppf})$.
Let $X, Y$ be algebraic spaces of finite type over $U$. By
Stacks, Lemma \ref{stacks-lemma-presheaf-mor-map-fibred-categories}
we obtain a map of presheaves
$$
\mathit{Mor}_{\Spacesstack_{ft}}(X, Y)
\longrightarrow
\mathit{Mor}_{\Spacesstack}(X, Y)
$$
which is an isomorphism as $\Spacesstack_{ft}$ is a full subcategory of
$\Spacesstack$. Hence the left hand side is a sheaf, because in
Lemma \ref{lemma-pre-stack-of-spaces}
we showed the right hand side is a sheaf. This proves (2).

\medskip\noindent
To prove condition (3) of
Stacks, Definition \ref{stacks-definition-stack}
we have to show the following: Given
\begin{enumerate}
\item a covering $\{U_i \to U\}_{i \in I}$ of $(\Sch/S)_{fppf}$,
\item for each $i \in I$ an algebraic space $X_i$ of finite type over $U_i$,
and
\item for each $i, j \in I$ an isomorphism
$\varphi_{ij} : X_i \times_U U_j \to U_i \times_U X_j$ of algebraic spaces
over $U_i \times_U U_j$ satisfying the cocycle condition over
$U_i \times_U U_j \times_U U_k$,
\end{enumerate}
there exists an algebraic space $X$ of finite type over $U$ and isomorphisms
$X_{U_i} \cong X_i$ over $U_i$ recovering the isomorphisms $\varphi_{ij}$.
This follows from
Bootstrap, Lemma \ref{bootstrap-lemma-descend-algebraic-space} part (2). By
Descent on Spaces, Lemma
\ref{spaces-descent-lemma-descending-property-locally-finite-presentation}
we see that $X \to U$ is of finite type which concludes the proof.
\end{proof}

\end{document}
